%% ABSTRACT
%% ============================================================================

Large Language Models (LLMs) have demonstrated remarkable capabilities in token-level pattern completion, matching or exceeding mid-level engineers on bounded coding tasks. However, when deployed in large codebases (50k+ lines of code), their quality collapses due to a fundamental structural blind spot: LLMs reason over token streams without access to the AST, call graph, or embedding manifold of the repository as a whole. This leads to failures in maintaining architectural invariants, understanding cross-file consequences, preserving cyclomatic complexity boundaries, and maintaining coherence with existing project conventions.

We present reasoning-core, a System 2 sidecar that augments LLM-based coding agents with structural reasoning capabilities. Our approach combines Tree-sitter parsing for AST and call graph extraction with a Mamba State Space Model (SSM) backbone (130M parameters) for semantic embedding. The sidecar computes an 8-dimensional risk vector capturing cyclomatic complexity, fan-in/out, depth, churn, coupling, cohesion, and novelty deltas between pre- and post-edit states.

In a historical controlled study across 8 real-world engineering tasks with 48
total runs (n=3 per task x setup), the reasoning-core sidecar reported:

- Historical token result: up to 29 percent reduction on cache-heavy tasks (29 percent on PR review, 27 percent on auth-abandonment)
- Historical token result: 8.2 percent average savings across all tasks
- Historical plan-quality difference: +0.32 (3.62 to 3.94 on 1-5 BARS scale)
- Historical implementation-quality difference: +0.20 (3.80 to 4.00)
- Historical test result: 100 percent pass rate on locked and rotated tests (vs 92 percent/90 percent for vanilla Claude Code)
- Historical diff-discipline difference: +0.23 (4.09 to 4.31)
- Historical repo-fit difference: +0.43 (3.71 to 4.14)

The system operates 100 percent locally with no external telemetry, processing
code entirely on the developer's machine. The study also measured an
approximately 98-second per-run overhead. These results are historical,
single-codebase evidence using an earlier risk-label schema; they do not
establish current real-usage quality, regression reduction, or net savings.
Current real-session validation is being conducted from labeled repository
sessions rather than a new synthetic pilot.

%% PLACEHOLDER: Add statistical significance results from iter-3 evaluation
%% PLACEHOLDER: Add comparison to other state-of-the-art approaches
%% PLACEHOLDER: Add theoretical contributions summary
